\chapter{把常见疑问串起来}
\section{我没有微调，VLM为什么能输出技能名}
模型可以在本次输入里读到技能清单及少量示例，然后按这些名称生成文字。这叫利用上下文执行指令，不需要当场更新权重。就像你给一个会读说明书的人一张按钮表，他可能按照说明选择按钮；是否能在复杂状态下一直选对，还需要实际验证。

本地微调把许多“这个输入应该选择这个技能”的对应关系用于更新参数，让模型更适应这套任务。技能库本身由代码定义。未来新增技能时，需要增加可执行控制器、参数约束、接口说明、示范和评测，不能只在提示词里多写一个不存在的动作名。

\section{颜色、位置和正确动作分别从哪里来}
“踢红球”里的目标颜色来自用户原话；延续性任务还可利用允许的上下文。当前两路RGB和RGB-D辅助测量提供可见球颜色与相对位置。规则老师据此选择动作，导出器把通过验收的选择写为目标标签。

模型输入中不会直接放入本次正确的 \texttt{skill}。评测专用的 \texttt{--expected-color} 也不进入模型输入。否则模型可能读取答案而不是理解目标，得到看似很高但无意义的分数。

\section{JSON格式是谁规定的，是否永远固定}
五字段格式由本项目的协议、提示词与解析器共同规定。Qwen本身是可生成文本的模型，不天然被限制为这五个字段。训练目标会强化这个格式；部署还应保留严格校验。改变字段名、允许的技能或时长范围后，应同步改变提示词、训练标签、解析器、执行器和测试。

当前离线生成评估不使用强制JSON语法解码，输出可能不合法。原有GGUF服务与客户端的约束方式则应按实际请求记录识别。比较两个后端时，不能一方强制格式而另一方无约束，却把全部差异归结为模型权重。

\section{SFT、LoRA和QLoRA是不是三种互斥训练}
它们描述不同维度。SFT说的是“拿示范输入和正确输出做监督训练”；LoRA说的是“只训练插入的小矩阵”；QLoRA进一步说的是“以低比特量化基础模型，训练LoRA”。因此本轮可以同时被称为SFT、QLoRA；其中动作示范部分也属于行为克隆。

同一个SFT目标，可以选择全量参数、部分层、LoRA或冻结特征后的预测头。DPO等则改变训练数据结构和目标，不能仅通过把优化器名字替换一下就变成另一种学习方法。

\section{预测技能概率和生成技能JSON有什么不同}
技能预测头可以一次输出九个技能的分数，再通过softmax得到归一化概率；另设任务类型、时长和颜色头，最后由程序组装JSON。它的输出类别由头的维度和标签表规定。

当前QLoRA方案逐token生成JSON。它在每个生成位置对词表中的token分配概率，经过多个token才形成 \texttt{move\_forward} 及其他字段。两个方案都可以学习技能选择，但损失形式、解码方式和格式错误类型不同。生成模型中的一个token概率，不能直接当成某个完整技能的可靠成功概率。

\section{动作是离散的，为什么还要预测时长}
技能名称是离散选择。前进0.10秒与前进0.60秒属于同一技能，但会造成不同位移和下一幅画面。持续时间决定一次决策的幅度，也影响闭环频率。

当前接口把时长限定为0.02秒的整数倍，转向和站立上限1秒，平移上限0.6秒，踢球固定0.5秒。因此可以把时长看作受限制的参数，也可以在新预测头方案里离散成类别。不能把“所有技能都离散”理解为每个技能永远只允许同一个执行时长。

\section{为什么一次只执行一个技能}
当前画面可能只能证明“向左转一点合理”，不能保证转完后前进再踢球仍然合理。短动作后再观察，使系统有机会纠正误差。一次生成整串动作会积累位置误差，还容易错过球或误碰其它球。

这不意味着模型每个物理步都要重新推理。一个短技能由底层控制器连续执行，上层在技能完成并恢复站立后再获取新观测。上层和底层分工，兼顾语言决策与稳定运动。

\section{训练是不是同时运行四只鸭子}
本轮不是在线与仿真交互的强化学习。数据采集阶段曾用三个独立仿真进程加快采集；正式训练阶段从磁盘读取样本，一个GPU训练进程每次处理一条样本，累计四条的梯度后更新一次。

GPU内部矩阵运算本来就有大规模并行，但这与同时放入几条样本、运行几个仿真环境是不同层面的“并行”。把梯度累计4改成8，只改变累计多少样本后更新，不会自动把读取和前反向变成八条同时处理。

\section{为什么只训练少量参数，仍然要十多个小时}
参数更新数量只影响部分成本。图像编码、长上下文语言主干前向、用于训练各层LoRA的梯度传播、梯度检查点重算和量化计算，都仍需执行。batch为1也限制了摊薄某些开销的机会。

冻结后缓存特征的方案能够将主干成本集中到首次特征提取，随后只训练小头，所以后续训练通常更轻。但在线运行新图像时仍要计算特征，不能直接用“缓存后训练很快”推断机器人每次决策同样很快。

\section{冻结、eval和no\_grad分别是什么意思}
\textbf{冻结参数}通常表示不为这些参数计算或保存需要更新的梯度。\textbf{eval模式}改变dropout等模块的运行行为，使推理更稳定，并不自动禁止梯度。\textbf{no\_grad}表示该计算范围不建立用于反向传播的梯度记录。

训练LoRA时，基础权重虽冻结，相关层仍需参与使LoRA可学习的计算图；不能随意把整个主干套进no\_grad，否则可能切断LoRA所需梯度。完全冻结特征提取器、只训练外接头时，则可以把特征提取放在no\_grad下。缓存特征时还应使用固定推理设置，避免dropout和输入处理不一致。

\section{成功轨迹里也会有多余动作，怎么办}
最终任务通过，只能证明整段轨迹达成了验收条件，不证明每一步都是最优。老师可能走了绕路或用了更多小调整。行为克隆会学习这些动作习惯，因此示范质量与覆盖率同样重要。

当前剔除了固定执行器恢复动作作为训练标签，但保留规则决策中的必要站稳。后续可以对成功轨迹进一步审查、缩短冗余动作；同时要保留真实可执行和稳定性证据，不能为了让轨迹更短而直接删除必要观察或站稳步骤。

\section{为什么失败轨迹没有直接加进训练}
失败轨迹的某些早期动作可能合理，也可能已经偏离目标。当前没有可靠的逐步纠正标注，因此保守地排除整个失败踢球任务的动作监督。这样避免把“最终失败的动作序列”作为应该照做的答案。

同一episode的独立视野问答可以是正确的，所以仍保留。将来若能给失败状态补上正确动作，就可以把这些纠正后的样本加入训练；这与把失败轨迹原样当成正确标签是两种操作。

\section{有RGB-D辅助，还算VLM吗}
模型仍然接收图像和语言，并利用这些信息进行决策；增加结构化测量并不改变它采用VLM的事实。但能力归因必须准确：当前球颜色分割和深度拟合由感知模块辅助，不能把所有几何检测成绩都称为VLM从RGB独立学到的能力。

若希望纯RGB决策，应设计单独的无检测输入实验，并重新检查距离判断、对齐和踢球表现。把已经训练依赖测量的模型直接删掉测量字段，只能算分布变化测试，不能预先保证性能相同。

\section{是不是需要VQ-VAE动作码本}
当前九个技能已有明确编号和执行器，因此可以直接预测技能类别或生成技能名称，不必再训练一个压缩连续动作轨迹的码本。VQ-VAE等方法适合某些希望从复杂动作序列中学习离散表示的方案，但不构成本项目技能分类的前置条件。

\section{“看到了什么”的功能怎样保留}
当前JSON接口包含 \texttt{answer} 和自然语言 \texttt{reply}，训练数据里也有视野问答。QLoRA能够把这类问答纳入同一个生成目标；但已有示范主要围绕三色球，开放环境描述仍需单独数据和评测。

冻结VLM加技能头时，可以保留原VLM语言生成分支来回答问题，由任务类型判断选择技能头还是语言回复。技能头自身只输出有限类别，不能独自生成任意场景描述。也需要防止问答请求误触发运动。

\section{自由表达是否需要每一句都写规则}
预训练语言能力可以让模型理解未逐字出现过的表达，但它能否把这些表达正确映射到你的技能，需要验证。数据应包含同一意图的多种说法，也需要同一画面对应不同目标和不同运动指令，避免模型只看图或记住固定任务顺序。

目前采集包含多种预先设计的表达，并通过独立复位运动数据减轻顺序捷径。它不等于已经覆盖所有口语、否定、组合任务或含糊指代。对“别踢红球，去踢蓝球”“先看看再往后退”等复杂表达，必须明确是否属于支持范围再测试。

\section{训练完是否还要API，是否收费}
当前训练和推理使用本地权重，不要求云端大模型API，因此没有这条流程产生的云API调用费；本地显卡、存储和电力仍有成本。官方模型卡声明Apache-2.0许可，使用时应按实际许可证要求处理。\cite{qwen-card}

将来可选择把部分对话或视觉判断放到远端服务，但这会改变延迟、费用、数据传输和故障方式，也需要验证接口与原流程的一致性。远端API是部署选择，不是完成这次本地技能微调的必要条件。

\section{是否应该现在就训练最简单的方案一}
可以把它作为独立对照：使用同一预训练VLM、同一输入与数据划分，提取不含目标答案的特征，训练小型技能头，再运行同样的闭环测试。当前尚未实现这条训练和部署入口，不能只改一个QLoRA参数来切换。

先完成正在运行的实验，有助于保留一个可比较的基准，也能腾出显存。小头通常训练更轻，但当基础特征对精确距离、历史决策或颜色指令区分不足时，学习上限也可能受限。最后选择哪条路线，应看真实任务与资源指标。

\section{读完后可以用这六个问题自查}
\begin{enumerate}
\item 能否指出模型当前看到哪些信息，以及哪些信息只用于验收？
\item 能否解释一个任务、一个样本、一次微批次与一次参数更新的差别？
\item 能否说清楚本轮loss计算的是哪些token，为什么四次反向只更新一次？
\item 能否区分模型输出合法、匹配示范、实际踢球成功三个层次？
\item 能否说明当前GGUF、基础safetensors、LoRA适配器和可恢复checkpoint分别做什么？
\item 能否为一个失败案例判断更可能需要补数据、改状态表示、改执行器，还是继续训练？
\end{enumerate}
如果其中某项还不清楚，回到对应章节，用真实样本沿着“画面和指令—输入消息—模型输出—执行器—新观测”的链路再走一遍，比单独记术语更有效。
